\gddchapter{Testing sequence}
The prototype can be played with both voice and keyboard and the navigation modes can be swapped with a button press. 


\section{Novel version B - Voice navigation}

Since the feedback from the last 3 tests the testing is now done in reverse order - the participant is playing the voice version first and then switching to the keyboard baseline version.


\subsection{Observation log}
\begin{table}[H]
\centering
\begin{tabular}{|p{5cm}|p{10cm}|}
\hline
\textbf{Field} & \textbf{Response} \\ \hline
\textbf{Physical comfort} How is the participant feeling & The participant seems to be comfortable \\[1.5cm] \hline
\textbf{Natural language observation} (is the player using long sentences or simple commands? What type of language do they use? Are they speaking first or 3rd person?) &The player is using simple voice commands in 1st person\\[1.5cm] \hline
\textbf{Linguistic guessing} (is the particpant trying to guess the possible commands by speaking to the game?) &The particpant is sticking to simple commands \\[1.5cm] \hline
\textbf{Processing time user reaction} (Does the wait time break the suspension of disbelief) & It's hard to observe since there doesn't seem to be a suspension of disbelief in the first place.\\[1.5cm] \hline
\textbf{Voice processing pipeline edge cases} (System edge cases eg. two vs to)& No drastic errors observed.\\[1.5cm] \hline
\end{tabular}
\end{table}


\section{Baseline A - Keyboard navigation}
After playing the audio version of the game, the player is now experiencing the baseline version with keyboard usage.
This allows the player to anchor in the typical experience and serves as a comparative point in the A/B testing.

\subsection{Observation log}
\begin{table}[h]
\centering
\begin{tabular}{|p{5cm}|p{10cm}|}
\hline
\textbf{Field} & \textbf{Response} \\ \hline
\textbf{Physical comfort} How is the participant feeling & The participant looks equally as comfortable as in the voice-driven version. \\[1.5cm] \hline
\textbf{Immersion markers} (flow \& frustration) & The player seems non immersed in both versions, focued more on the novelty of the fact that she's in an experiment.\\[1.5cm] \hline
\end{tabular}
\end{table}


\subsection{Sticky moments}
The moments that stood out during testing marked with timestamps.

The paper document with the list of the sticky moments got lost for this participant,
but going from the researcher memory there doesn't seem to be any new important notes that would differ from the previous sessions.
The system failed on the lounge again and the participant was confused with the UX of recording. On top of that the stair navigation was proven
to be difficult, per usual.

% \begin{table}[h]
% \centering
% \begin{tabular}{|p{3cm}|p{10cm}|}
% \hline
% \textbf{Timestamp} & \textbf{Log} \\ \hline
% 05:00 & Start of the playtest \\ \hline 
% 06:00 & Long descriptions \\ \hline 
% 10:45 & Not very obvious how to use items \\ \hline 
% 09:10 & "I could skip reading this. Too much time spent on reading" \\ \hline 
% 10:00 & Stair name confuision \\ \hline 
% 12:00 & Issues with UX - recorning \\ \hline 
% 13:00 & Double VUI failure \\ \hline 
% 16:00 & Manual interverion lounge \\ \hline 
% 16:00 & Pharmacy 1 time error than worked \\ \hline 
% \end{tabular}
% \end{table}

\section{Alignment check}
How much was the researcher participating in the gameplay experience? 

The researcher did his best not to participate actively during the testing sequence. 



